\section{Ranks in a scalar balanced-compound family}
\label{sec:balanced-compound}

The two-by-two block structure of the highest normalizer coefficient
suggests examining a scalar compound as a crossing deformation in its
own right. This family admits an exact rank decomposition into
intersection matrices. Throughout this section the field has
characteristic two.

Fix $n\geq1$, set $M=2n$ and $q=\binom{2n}{n}$, and put
\[
 L(z)=\begin{pmatrix}I_M&zI_M\\ I_M&I_M\end{pmatrix}.
\]
Index each terminal by an $n$-subset of $[M]$. For indices $A,B,C,D$,
let $T(z)_{A,B,C,D}$ be the determinant of the $M$ by $M$ minor
whose rows are $A$ in the first block and $C$ in the second, and
whose columns are $D$ in the first block and $B$ in the second.

\begin{proposition}[The scalar compound and its rank blocks]
\label{prop:scalar-compound-blocks}
The tensor $T(z)$ has degree at most $n$, and
\[
 T(1)_{A,B,C,D}=[C=A^{\mathrm c}][B=D^{\mathrm c}].
\]
Thus $T(1)$ is a pair of dimension-$q$ crossing wires after local
complement permutations. Its reflection is
\begin{equation}\label{eq:scalar-compound-reflection}
 z^nT(1/z;A,B,C,D)=T(z;A,D,C,B).
\end{equation}
Write $T(z)=\sum_{j=0}^n z^jT_j$ and
$T_a=\sum_{j=0}^n a_jT_j$. Define
\[
 d_k=\binom{2n}{k}\binom{2n-k}{k}\qquad(0\leq k\leq n).
\]
For $r=n-k$, let $H_k(a)$ be the matrix indexed by $r$-subsets
$S,U$ of a $2r$-element set, with entries
\[
 H_k(a)_{S,U}=\sum_{\ell=0}^{r}
                a_{k+\ell}\binom{|S\cap U|}{\ell}.
\]
Then the exact adjacent-cut rank is
\begin{equation}\label{eq:scalar-compound-rank}
 \operatorname{rank}_{ad\mid bc}T_a
 =\sum_{k=0}^n d_k\operatorname{rank}H_k(a).
\end{equation}
\end{proposition}
\begin{proof}
The one-particle matrix couples only equal mode indices. A minor can
therefore be nonzero only if each mode has the same row and column
multiplicity. Equivalently,
\[
 P=A\setminus D=B\setminus C,\qquad
 Q=D\setminus A=C\setminus B.
\]
These disjoint sets have the same size $k$. A mode in $P$ contributes
$z$, a mode in $Q$ contributes one, and a mode occupied twice in
both the row and column sets contributes $1+z$. All other nonzero
one-mode factors are one. If $v=|A\cap C|$, the entry is therefore
\[
 z^k(1+z)^v.
\]
Since $P$ and $A\cap C$ are disjoint subsets of $A$, we have
$k+v\leq n$, proving the degree bound. At $z=1$ only $v=0$
survives, which forces $C=A^{\mathrm c}$ and $B=D^{\mathrm c}$.
Interchanging $B,D$ replaces $k$ by $n-v-k$, proving the reflection.

In the $ad\mid bc$ flattening, distinct ordered pairs $(P,Q)$
index disjoint row and column blocks. There are $d_k$ such pairs
of size $k$. Inside one block write
\[
 A=P\cup S,\quad D=Q\cup S,\quad
 B=P\cup U,\quad C=Q\cup U.
\]
The two common subsets $S,U$ have size $n-k$ in the complement
of $P\cup Q$, which has size $2(n-k)$. The block entry is
$z^k(1+z)^{|S\cap U|}$. Extracting the chosen linear combination
of coefficients gives $H_k(a)$, and additivity of matrix rank over
disjoint blocks proves \eqref{eq:scalar-compound-rank}.
\end{proof}

\begin{theorem}[The scalar coefficient span has exponent at least three halves]
\label{thm:scalar-compound-barrier}
For arbitrary coefficients in any characteristic-two extension
field, let $A=T_a$ and $B=T(1)-A$. For every $k$,
\[
 \operatorname{rank}_{ad\mid bc}A+
 \operatorname{rank}_{ab\mid cd}B
 \geq\min(d_k,d_{n-k}).
\]
In particular this sum is at least
\[
 \frac{q^{3/2}}{2(n+1)}.
\]
Consequently the total rank of these decompositions has asymptotic
exponent at least $3/2$ as $n\to\infty$.
\end{theorem}
\begin{proof}
The $t=1$ specialization of
\eqref{eq:scalar-compound-word-bound}, proved below from the
projected pair-polynomial theorem, is a sum of the nonnegative
terms $\min(d_k,d_{n-k})$. This proves the first bound.

Take $k=\lfloor n/2\rfloor$. If $n$ is even,
\[
 d_k=q\binom{n}{n/2}\geq\frac{q2^n}{n+1}.
\]
If $n$ is odd, the smaller of $d_k,d_{n-k}$ is
\[
 \frac q2\binom{n}{\lfloor n/2\rfloor}
 \geq\frac{q2^{n-1}}{n+1}.
\]
The central binomial coefficient is at least the average of all
$n+1$ coefficients. Finally $q\leq4^n$, so $2^n\geq\sqrt q$.
This proves the displayed uniform bound and its exponent.
\end{proof}

\begin{theorem}[Tensor powers of the full scalar coefficient span]
\label{thm:scalar-compound-tensor-barrier}
Fix $n$. For $t\geq1$, let $A$ be any element of the span of
\[
 T_{j_1}\otimes\cdots\otimes T_{j_t},\qquad 0\leq j_i\leq n,
\]
with the four terminal groups retained, and put $B=T(1)^{\otimes t}-A$.
The sum of the $ad\mid bc$ rank of $A$ and the $ab\mid cd$ rank of
$B$ is at least
\begin{equation}\label{eq:scalar-compound-word-bound}
 \sum_{\mathbf j\in\{0,\ldots,n\}^t}
 \min\left\{\prod_{i=1}^t d_{j_i},\,
                  \prod_{i=1}^t d_{n-j_i}\right\}.
\end{equation}
The exponential rate of this lower bound is
\[
 D_n=\sum_{k=0}^n\sqrt{d_kd_{n-k}}\geq q^{3/2}.
\]
Thus every such sequence of decompositions has crossing exponent
at least $3/2$, including arbitrary coefficients over
characteristic-two extension fields.
\end{theorem}
\begin{proof}
Apply Theorem~\ref{thm:projected-pair-barrier} to the coordinate set
of weight-$n$ strings on $2n$ sites, using the exact coefficient
identification and transpose transport in
Proposition~\ref{prop:balanced-pair-identification}.
For a subset $A$ of size $2j$ the two weights are
\[
 w(A)=\binom{2n-2j}{n-j},\qquad
 w(A^{\mathrm c})=\binom{2j}{j};
\]
odd cardinalities have weight zero. There are $\binom{2n}{2j}$
subsets of size $2j$. Multiplying this common multiplicity into
each alternative turns the two weights into $d_{n-j}$ and $d_j$.
On $t$ copies this gives exactly
\eqref{eq:scalar-compound-word-bound}, allowing arbitrary
coefficient vectors. The fixed-coordinate-set rate in
\eqref{eq:projected-pair-rate} becomes
$D_n=\sum_j\sqrt{d_jd_{n-j}}\geq q^{3/2}$.
\end{proof}

The block formula also makes small exact searches inexpensive. The
following minima range over all $2^{n+1}$ coefficient vectors
$a\in\mathbb F_2^{n+1}$; the listed rank pair is an attaining choice.
\begin{center}
\begin{tabular}{@{}rrrrr@{}}
\toprule
$n$ & $q$ & First cut & Second cut & Minimum total\\
\midrule
$1$&$2$&$1$&$1$&$2$\\
$2$&$6$&$1$&$15$&$16$\\
$3$&$20$&$36$&$36$&$72$\\
$4$&$70$&$63$&$783$&$846$\\
$5$&$252$&$1926$&$1926$&$3852$\\
$6$&$924$&$4356$&$44188$&$48544$\\
\bottomrule
\end{tabular}
\end{center}
These finite minima concern coefficients over $\mathbb F_2$;
Theorem~\ref{thm:scalar-compound-barrier} holds over every extension
without that search restriction. The determinant definition and
the block reduction were also compared directly in the first two
cases, using respectively $16$ and $1296$ minors.

The parameter-dependent exterior changes of basis connecting the
digit kernel to a scalar block matrix need not preserve coefficient
ranks. Thus this obstruction leaves the coefficient-rank problem
for the original digit construction unresolved.
